\documentclass[12pt,english]{article}
\usepackage[T1]{fontenc}
\usepackage[latin9]{inputenc}
\usepackage[toc,page]{appendix}
\usepackage{geometry}
\geometry{verbose,tmargin=1.5in,bmargin=1.5in,lmargin=1.5in,rmargin=1.5in}
\usepackage{bm}
\usepackage{amstext}
\usepackage{comment}
\usepackage{amsmath}
\usepackage{hyperref}
\hypersetup{
     colorlinks = true,
     citecolor = navy,
     linkcolor = navy,
     urlcolor=navy
}
\usepackage{adjustbox}
\usepackage[space]{grffile}
\usepackage{setspace}
\usepackage[normalem]{ulem}
\usepackage{nicematrix}
\usepackage{xcolor}
\usepackage{soul}
\usepackage{arydshln}
\usepackage{dsfont}
\usepackage{refstyle}
\usepackage[capitalize,nameinlink]{cleveref}
\usepackage{setspace}
\usepackage{bbm}
\usepackage[tableposition=t]{caption} % ensures caption of longtables is normalsized.
\usepackage{mathpazo} % add Palatino font
\usepackage{longtable}
\usepackage{amssymb}
\usepackage{amsfonts}
\usepackage{amsmath}
%\usepackage[margin=1in]{geometry} % to set margins
\usepackage{scalefnt}
\usepackage{wrapfig}
\usepackage{graphicx}
\usepackage{graphics}

%\usepackage{subfigure}
\usepackage{lscape}
\usepackage{pdflscape}
\usepackage{multirow}
\usepackage{array}
\usepackage{rotating}
\usepackage{bibunits}
\usepackage{fix-cm}
\usepackage{float}
\floatplacement{table}{H}
\usepackage{epsfig}
\usepackage{epstopdf}
\usepackage{refstyle}
\usepackage{tikz}
\usepackage{booktabs}
\usepackage{tabularx}
\usepackage{array}
\newcolumntype{x}[1]{>{\raggedleft\arraybackslash}p{#1}}%
\newcolumntype{d}[1]{>{\centering\arraybackslash}p{#1}}%
\newcolumntype{y}[1]{>{\raggedright\arraybackslash}p{#1}}%
\usepackage{verbatim}
%\epstopdfsetup{suffix=}
\setcounter{MaxMatrixCols}{10}
\usepackage[longnamesfirst,authoryear]{natbib}
\newenvironment{proof}[1][Proof]{\noindent \textbf{#1.}}{\rule{0.5em}{0.5em}}
\setlength{\topmargin}{-0.80in} \setlength{\textheight}{9.50in}
\setlength{\evensidemargin}{.0in} \setlength{\oddsidemargin}{.0in}
\setlength{\textwidth}{6.85in}
\setlength{\voffset}{0.4cm}
\setlength{\hoffset}{-0.5cm}
\usepackage{setspace}
\usepackage{afterpage}
\usepackage[flushleft]{threeparttable}
\usepackage{tabulary}
%\usepackage[titletoc,title]{appendix}
\usepackage[bottom]{footmisc}
\usepackage{color}
	\definecolor{maroon}{RGB}{128,0,0}
	\definecolor{navy}{RGB}{0,0,165}
	\definecolor{lb}{RGB}{150,165,190}
	\definecolor{bluenice}{rgb}{0.25,0.51,0.76}
	\definecolor{rednice}{rgb}{0.85,0.26,0}
	\definecolor{orangenice}{HTML}{EB811B}
	\definecolor{greennice}{HTML}{14B03D}
\usepackage{colortbl}
\usepackage{caption}
\usepackage{subcaption}
\usepackage{placeins}
\usepackage{nicefrac}
%\usepackage{chngcntr}
\newtheorem{assumption}{Assumption}
\newtheorem{definition}{Definition}
\usetikzlibrary{matrix,decorations.pathreplacing,arrows.meta,positioning}

\newcommand{\vy}{\underline{y}}
\newcommand{\ve}{\underline{\epsilon}}
\newcommand{\A}[1]{\mathbold{A_{#1}}}
\newcommand{\D}[1]{\mathbold{D_{#1}}}
\newcommand{\C}[1]{\mathbold{C_{#1}}}
\newcommand{\B}[1]{\mathbold{B_{#1}}}
\newcommand{\R}[1]{\mathbold{R_{#1}}}

%\newtheorem{definition}{Definition}
\begin{document}
\title{Original Inflation Calculation with Leontief}
\author{N}
\maketitle

\paragraph{Purpose: }

I want to derive an approximation that will connect the Leontief ShapirIO decomposition to net inflation defined the way Shapiro does it.

Let's get started.

\section{Derivations*}

CPI inflation is now calculated as:

\[
	\pi_{CPI,t} \equiv \sum_p \omega_{PCE,p,t} \underbrace{\left(\frac{P_{p,t}}{P_{p,t-1}} - 1\right)}_{\equiv \pi_{p,t}} = \underline{\omega_{PCE,t}}'\underline{\pi}_{t}
\]

Note that our I-O decomposition has the following decomposition

\[
	\underline{P}_t = \underbrace{(I-A_t)^{-1}diag(\mathds{1}_p'A_t)}_{\equiv B_t} \underline{P}_{VA,t}
\]
where I put a $t$ subscript on $A_t$ because we have changes in the set of available products.

In summation terms this is:

\[
	P_{p,t} = \sum_{p'} B_{p,p',t} P_{VA,p',t}
\]

Now we can define product level inflation as:

\[
	\pi_{p,t} \equiv \frac{P_{p,t}-P_{p,t-1}}{P_{p,t-1}} = \frac{\sum_{p'}\left(B_{p,p',t}(P_{VA,p',t} - P_{VA,p',t-1}) + (B_{p,p',t} - B_{p,p',t-1})P_{VA,p,t-1}\right)}{\sum_{p''} B_{p,p'',t-1}P_{VA,p'',t-1}}
\]

where in the numerator I added and subtracted $B_{p,p',t}P_{VA,p,t-1}$ to get this expression. Next we can do the following adjustments as well:

\begin{align*}
 \pi_{p,t} &= \frac{\sum_{p'}\left(B_{p,p',t}(P_{VA,p',t} - P_{VA,p',t-1})\textcolor{bluenice}{\frac{P_{VA,p',t-1}}{P_{VA,p',t-1}}} + (B_{p,p',t} - B_{p,p',t-1})P_{VA,p,t-1}\textcolor{orangenice}{\frac{B_{p,p',t}}{B_{p,p',t}}}\right)}{\sum_{p''} B_{p,p'',t-1}P_{VA,p'',t-1}}\\
 &=\frac{\sum_{p'}\left(B_{p,p',t}\textcolor{bluenice}{P_{VA,p',t-1}}\frac{P_{VA,p',t} - P_{VA,p',t-1}}{\textcolor{bluenice}{P_{VA,p',t-1}}} +\textcolor{orangenice}{B_{p,p',t}}P_{VA,p,t-1} \frac{B_{p,p',t} - B_{p,p',t-1}}{\textcolor{orangenice}{B_{p,p',t}}}\right)}{\sum_{p''} B_{p,p'',t-1}P_{VA,p'',t-1}}\\
 &=\underbrace{\sum_{p'}\frac{B_{p,p',t}P_{VA,p',t-1}}{\sum_{p''} B_{p,p'',t-1}P_{VA,p'',t-1}}}_{\equiv \omega_{VA,p,p',t}}\left(\pi_{VA,p',t-1} + \Delta B_{p,p',t}\right)
\end{align*}

where $\pi_{VA,p',t}\equiv \frac{P_{VA,p',t} - P_{VA,p',t-1}}{P_{VA,p',t-1}}$ and $\Delta B_{p,p',t} \equiv \frac{B_{p,p',t} - B_{p,p',t-1}}{\textcolor{orangenice}{B_{p,p',t}}}$

So far this is an exact decomposition \textcolor{orange}{(don't forget the time subscript in the change in $B$ matrix -- we divide by $B$ in time $t$ not $t-1$.}

Now let's link this to CPI inflation:

\begin{align*}
	\pi_{CPI,t} &= \sum_p \omega_{PCE,p,t}\pi_{p,t}\\
	& = \sum_p \omega_{PCE,p,t}\sum_{p'}\omega_{VA,p,p',t} \left(\pi_{VA,p',t-1} + \Delta B_{p,p',t}\right)\\
	& = \sum_{p'}\sum_p \omega_{PCE,p,t}\omega_{VA,p,p',t} \left(\pi_{VA,p',t-1} + \Delta B_{p,p',t}\right)\\
	&=\sum_{p'}\left(\underbrace{\sum_p \omega_{PCE,p,t}\omega_{VA,p,p',t}}_{\equiv W_{p',t}} \pi_{VA,p',t-1} + \underbrace{\sum_p \omega_{PCE,p,t}\omega_{VA,p,p',t}\Delta B_{p,p',t}}_{\equiv V_{p',t}}\right)
	&= \sum_{p'}W_{p',t}\pi_{VA,p',t-1} + V_{p',t}
\end{align*}

And this is also an exact decomposition!

\section*{Implementation Steps}

\begin{enumerate}
	\item Start by constructing the time varying matrices $B_t\equiv (I-A_t)^{-1}diag(\mathds{1}_p'A_t)$ which should be a different $\mathcal{P}_t$ x $\mathcal{P}_t$ for each time period (where $\mathcal{P}_t$ is the number of products in that period). Note also that this only differs by the set of products -- if that set doesn't change, then $B_t = B_{t-1}$.
	\item Then use this to construct $\omega_{VA,p,p',t} \equiv \frac{B_{p,p',t}P_{VA,p',t-1}}{\sum_{p''} B_{p,p'',t}P_{VA,p'',t-1}}$. Note that this is a different $\mathcal{P}_t$ x $\mathcal{P}_t$ matrix for each time period $t$. 
	\item Calculate value added inflation and the change in $B$ matrices (not the denominator in the $B$ formula):
	\begin{align*}
		\pi_{VA,p',t}&\equiv \frac{P_{VA,p',t}-P_{VA,p',t-1}}{P_{VA,p',t-1}}\\
		\Delta B_{p,p',t} &\equiv \frac{\Delta B_{p,p',t}-]\Delta B_{p,p',t-1}}{\Delta B_{p,p',t}}
	\end{align*}
	\item Then create $W_{p',t}$ and $V_{p',t}$
	
	\begin{align*}
		W_{p',t}&\equiv \sum_p \omega_{PCE,p,t}\omega_{VA,p,p',t}\\
		V_{p',t}&\equiv \sum_p \omega_{PCE,p,t}\omega_{VA,p,p',t}\Delta B_{p,p',t}
	\end{align*}
	\item Now with these check that the following holds:
	\[
	\pi_{CPI,t} = \sum_{p'} \left(W_{p',t}' \pi_{VA,p',t} + V_{p',t}\right)
	\]
	\item If not, let Nick know!
	\item Then calculate demand and supply inflation as follows:
	\begin{align*}
		\pi_{demand,t}^{Shapiro} &\equiv \sum_{p'} D_{p',t} \omega_{PCE,p',t} \pi_{p',t}\\
		\pi_{supply,t}^{Shapiro} &\equiv \sum_{p'} (1-D_{p',t}) \omega_{PCE,p',t} \pi_{p',t}\\		
		\pi_{demand,t}^{ShapirIO} &\equiv \sum_{p'} D_{p',t} \left(W_{p',t}' \pi_{VA,p',t} + V_{p',t}\right)\\
		\pi_{supply,t}^{ShapirIO} &\equiv \sum_{p'} (1-D_{p',t}) \left(W_{p',t}' \pi_{VA,p',t} + V_{p',t}\right)\\	
	\end{align*}
	\item Do the 12 month averaging thing you already do and then the standard plots.
\end{enumerate}

\end{document}